\documentclass[11pt]{article}

\usepackage{amsmath}
\usepackage{booktabs}
\usepackage[most]{tcolorbox}

\newtcolorbox{eqbox}{colback=blue!6!white, colframe=blue!45!black,
  boxrule=0.8pt, sharp corners, left=2pt, right=2pt, top=2pt, bottom=2pt}
\newtcolorbox{whybox}{colback=black!4!white, colframe=black!4!white,
  borderline west={2pt}{0pt}{black!30}, sharp corners,
  left=6pt, right=6pt, top=4pt, bottom=4pt}

\title{Gamma-Conditioned Bands on a Single-Name Equity (SPCX)}
\author{Alex Marthan}
\date{July 2026}

\begin{document}
\maketitle

\section{Setup and Benchmarks}

This is a companion to \emph{Backtest Results and Possible Improvements}, which studied
static and dynamic bands on ES. Here the same band machinery is applied to \textbf{SPCX},
a Nasdaq-listed single-name equity, using the SPCX data pipeline and cost model built on
the \texttt{Timmy/strategy-comparison} branch (notebook
\texttt{05\_spcx\_extension}); the band shapes and the analysis below are from notebook
\texttt{06\_spcx\_dynamic\_bands}. The sample is 22 trading days
(June 16 -- July 17, 2026), 120{,}951 simulated SPCX option executions, 1{,}716
five-minute bars (78 per day), with book delta \emph{and} gamma repriced by
Black--Scholes at every bar.

Two things differ materially from the ES study and drive every number below.

\textbf{Units.} One option contract delivers 100 shares and the hedge instrument is the
stock itself (\$1 of P\&L per share per \$1 move), so book delta is measured in
\emph{shares} and dollar P\&L needs no contract multiplier. Delta risk is the RMS of
realized residual-delta dollar P\&L per bar (residual shares $\times$ \$ price move).

\textbf{Cost model.} The ES impact level was anchored at ``$\approx0.5$bp per 1-lot at
median ADV,'' a figure that silently embeds ES's $\approx15\%$ volatility and its
negligible clip-to-ADV ratio. Neither holds here, so the SPCX cost model is rebuilt
rather than ported:
\begin{itemize}
\item \textbf{Square-root law, volatility-anchored.} The impact fraction of notional is
  $\alpha\sqrt{q/L}$ with $\alpha = Y\,\sigma_{\text{daily}}$ (Toth--Bouchaud prefactor
  $Y = 0.5$), giving $\alpha = 0.0273$ from SPCX's median ATM implied volatility.
\item \textbf{Execution-window liquidity.} Impact is charged against the liquidity
  available over the fill's window, $L = \text{ADV}\times 5/390$ for an intraday child,
  not against full-day ADV. This is what makes intraday hedging genuinely expensive here.
\item \textbf{Closing cross.} A fill on the final bar leans on full-day ADV, crosses at
  the auction price (zero spread), and takes a $0.30$ haircut on impact.
\end{itemize}

Over the sample SPCX traded \$123.12--\$222.80 (median \$155.27), with ATM implied
volatility of 70--126\% and a 20-day ADV proxy near 982{,}000 shares, about \$152M a
day. The measured median quoted half-spread is \$0.027 per share. The book is large
relative to that liquidity: the selected strategy turns over roughly 638{,}000 shares a
day, about 65\% of ADV, which is why impact rather than spread dominates.
Configurations are scored by the standardized loss
\[
L_{\text{std}} = \widetilde{C}_{\text{exec}} + \lambda\,\widetilde{C}_{\text{risk}},
\qquad \lambda = 1,
\]
normalizing cost by Full Hedge and risk by Zero Hedge.

\begin{table}[htbp]
\centering
\small
\begin{tabular}{lrrr}
\toprule
\textbf{Benchmark} & \textbf{Exec.\ cost} & \textbf{Risk (RMS/bar)} & $L_{\text{std}}$ \\
\midrule
Zero Hedge (flatten at close only) & \$9{,}974{,}886   & \$435{,}115 & 1.027 \\
Full Hedge (re-flatten every bar)  & \$367{,}974{,}631 & \$0         & 1.000 \\
\bottomrule
\end{tabular}
\caption{Benchmark bounds over the 22-day SPCX backtest. Full hedging a book this size
against one bar's worth of liquidity costs \$368M; never hedging carries \$435k of
residual-delta P\&L risk per bar.}
\end{table}

\section{The Gamma-Conditioned Band}

Only the gamma mechanism is carried over from the ES study. The time-decay factor
$\sqrt{R(t)/R(0)}$ is \emph{dropped}: its derivation assumed the forced end-of-day
flatten is expensive, and SPCX's Nasdaq closing cross makes that flatten cheap, removing
the penalty the shape was built to avoid. The band is therefore

\begin{eqbox}
\begin{equation}
\lvert \Delta(t) \rvert > H(t)(1+\epsilon),
\end{equation}
\begin{equation}
H(t) = \frac{H_0}{1 + c\,\lvert\Gamma(t)\rvert},
\end{equation}
\end{eqbox}

where $\Gamma(t)$ is the book gamma repriced each bar (shares of delta regenerated per
\$1 move) and $c \to 0$ recovers the static band. Note $c$ is \emph{not} dimensionless in
this form: it carries units of \$ per share, and the band halves at
$\lvert\Gamma\rvert = 1/c$. This book's median $\lvert\Gamma\rvert$ is 2{,}176 shares/\$, so
the useful range of $c$ is order $10^{-4}$--$10^{-3}$; the grid is
$c \in \{2,5,10,20\}\times10^{-4}$. (Dividing $\lvert\Gamma\rvert$ by its median and
rescaling $c$ gives the same family of bands with a dimensionless $c$.)
The rationale is that delta diffuses through the band at speed
$d\Delta \approx \Gamma\,dS$ (proposal Eq.~32): a high-gamma book breaches a fixed band
faster, so tightening when $\lvert\Gamma\rvert$ is high keeps breach dynamics comparable
across gamma regimes.

\textbf{Is there signal?} The band can only help if gamma varies intraday and carries
information not already in $\lvert\Delta\rvert$. Both hold, modestly: across bars
$\rho(\lvert\Delta\rvert,\lvert\Gamma\rvert) = 0.21$ (largely independent signal), the
within-day coefficient of variation of $\lvert\Gamma\rvert$ is $0.46$, and the median
day's intraday max/median gamma ratio is $2.1\times$. That $2.1\times$ is the whole range
over which the band can flex, and it caps how large the effect can be.

\section{Static vs.\ Gamma: Unconstrained Grid}

Both families sweep $H_0$, execution horizon and method at fixed $\epsilon = 5\%$ and
$T_{\text{cooldown}} = 15$ min; the gamma family adds the tilt $c \in \{2,5,10,20\}\times10^{-4}$. The
static grid is extended down to $H_0 = 5{,}000$ shares, well below the 50k floor used in
notebook \texttt{05\_spcx\_extension}.

\begin{table}[htbp]
\centering
\small
\begin{tabular}{lrrrr}
\toprule
\textbf{Strategy} & \textbf{Exec.\ cost} & \textbf{Risk/bar} & $L_{\text{std}}$ & \textbf{Hedges/day} \\
\midrule
Best static & \$88.5M & \$112{,}916 & 0.5001 & 47.6 \\
Best gamma  & \$88.2M & \$113{,}396 & 0.5004 & 46.5 \\
\bottomrule
\end{tabular}
\caption{Unconstrained grid winners: static $H_0 = 5{,}000$ shares, VWAP 30min; gamma
$H_0 = 25{,}000$ shares, $c = 2\times10^{-3}$, VWAP 30min. The band cuts the standardized objective by
$\approx51\%$ vs.\ Zero Hedge and $\approx50\%$ vs.\ Full Hedge -- the first-order
result. Static and gamma are within 0.1\% of each other.}
\end{table}

\begin{whybox}
\textbf{The width trap, again -- and a boundary solution.} The best unconstrained
configuration is the \emph{tightest static band in the grid}, firing 47.6 hedges/day.
This is an artifact, not a strategy. Impact is superlinear in clip size
($\propto q^{1.5}$), so splitting a parent hedge into $n$ children costs
$n(Q/n)^{1.5} = Q^{1.5}n^{-1/2}$ -- strictly falling in $n$ -- and the model charges no
fixed per-ticket cost to stop it. The static ladder confirms the degeneracy: $L$ falls
monotonically into the grid edge ($H_0 = 50$k: 0.5287; 25k: 0.5238; 10k: 0.5088;
5k: 0.5001), so the optimum sits \emph{on the boundary} and would keep moving if the grid
were extended. Any comparison of band \emph{shapes} at the unconstrained optimum is
therefore uninformative.
\end{whybox}

\section{Width-Matched Comparison}

Setting a static band to each gamma configuration's mean effective width (same horizon,
method, $\epsilon$, cooldown) isolates the value of \emph{conditioning on gamma} from the
value of simply being \emph{tighter}.

\begin{table}[htbp]
\centering
\small
\begin{tabular}{rrrrr}
\toprule
$H_0$ (sh) & Mean width (sh) & Gamma $L_{\text{std}}$ & Matched static $L_{\text{std}}$ & Edge \\
\midrule
25{,}000  & 9{,}049   & 0.5243 & 0.5248 & $+0.11\%$ \\
50{,}000  & 18{,}098  & 0.5268 & 0.5324 & $+1.06\%$ \\
100{,}000 & 36{,}195  & 0.5288 & 0.5386 & $+1.82\%$ \\
200{,}000 & 72{,}391  & 0.5292 & 0.5480 & $+3.42\%$ \\
300{,}000 & 108{,}586 & 0.5904 & 0.5389 & $-9.55\%$ \\
\bottomrule
\end{tabular}
\caption{Width-matched gamma vs.\ static at $c = 10^{-3}$, TWAP 60min. Positive = gamma better.}
\end{table}

The edge is positive across every band width that actually binds ($\approx$9k--72k shares)
and \emph{grows with width}, from $+0.11\%$ to $+3.42\%$ -- wider bands bind less often, so
the choice of \emph{when} to hedge matters more. It collapses to $-9.6\%$ at the widest
setting, where the band is so wide it rarely binds at all and the gamma tilt merely
over-tightens on the few occasions it does. So
gamma conditioning is genuinely risk-efficient width-for-width, but it is a second-order
refinement, not a free lunch -- consistent with the $2.1\times$ ceiling on intraday gamma
variation.

\section{The Operational Regime: A Hedge-Frequency Cap}

The unconstrained winner hedges 47.6 times a day, which no desk would do in this name.
Because the missing ingredient is a fixed per-ticket cost (Section~7), the cleanest
stand-in is to cap the average number of hedges per day and re-optimize both families
over the same grid.

\begin{table}[htbp]
\centering
\small
\begin{tabular}{lrrrl}
\toprule
\textbf{Cap/day} & \textbf{Static $L$} & \textbf{Gamma $L$} & \textbf{Edge} & \textbf{Winner} \\
\midrule
5    & 0.5242 & 0.5444 & $-3.86\%$ & static \\
10   & 0.5242 & 0.5338 & $-1.83\%$ & static \\
15   & 0.5242 & 0.5241 & $+0.01\%$ & level \\
20   & 0.5238 & 0.5118 & $+2.29\%$ & gamma \\
30   & 0.5238 & 0.5086 & $+2.91\%$ & gamma \\
\midrule
none & 0.5001 & 0.5004 & $-0.06\%$ & static \\
\bottomrule
\end{tabular}
\caption{Best static vs.\ best gamma under a cap on hedge frequency. Positive = gamma
better.}
\end{table}

This is the decision-relevant result. Below $\approx$10 hedges/day both families are
forced so wide that the band rarely binds and shape barely matters, so the simpler static
rule wins. The two draw level at $\approx$15/day, and from $\approx$20/day upward -- the
realistic operating range -- the gamma-conditioned band wins outright, by $+2.3\%$ and
$+2.9\%$ of standardized loss at caps of 20 and 30. The
mechanism is intuitive: a fixed hedge budget is precisely the constraint that makes
\emph{where} you spend hedges matter, and spending them at high convexity is what gamma
conditioning optimizes. The gamma advantage is thus a \emph{capacity-regime} effect, not
a $\lambda$ effect: unconstrained, the degenerate tight static band wins at every
$\lambda \in [0.25, 4]$.

\section{Band-Edge Targeting}

Replacing the full flatten with the Whalley--Wilmott-style edge target
$Q_t = -\,\mathrm{sign}(\Delta(t))\,(\lvert\Delta(t)\rvert - H(t))$ parks the residual at
the band edge rather than round-tripping through zero.

\begin{table}[htbp]
\centering
\small
\begin{tabular}{lrrrrr}
\toprule
$H_0$ (sh) & 10{,}000 & 25{,}000 & 50{,}000 & 100{,}000 & 200{,}000 \\
\midrule
Static band & $+0.59\%$ & $+1.14\%$ & $+2.00\%$ & $+0.31\%$ & $-3.22\%$ \\
Gamma band  & $+0.10\%$ & $+0.36\%$ & $-0.04\%$ & $-0.20\%$ & $-2.43\%$ \\
\bottomrule
\end{tabular}
\caption{Improvement from edge targeting (TWAP 60min). Positive = better than full
flatten.}
\end{table}

Edge targeting is markedly less valuable on SPCX than on ES, where it was the strongest
single refinement. The reason is the closing cross: carrying residual delta to a cheap
auction close is most of what edge targeting buys, and on SPCX the auction already
delivers that to every strategy. It remains mildly positive for moderate \emph{static}
widths and backfires for wide bands under $\lambda = 1$; stacked on the gamma band it adds
essentially nothing, turning negative beyond $25{,}000$ shares. The gamma band already
tightens where the residual would otherwise be parked, leaving edge targeting little to do.

\section{Cost-Model Discussion}

Both cost components are size-proportional -- spread (linear in shares) and square-root-law
impact (superlinear, $\propto \lvert q \rvert^{1.5}$ once the $q\sqrt{q/L}$ form is
expanded) -- so both vanish as orders shrink and the model has no mechanism to penalize
frequent triggering. This is the same structural gap identified in the ES study, and on
SPCX it is more consequential: charging impact against one bar's liquidity makes each
individual hedge expensive, which pushes the optimizer toward many small clips even
harder. The missing term is a \textbf{fixed per-ticket cost}: a charge per child order
independent of size (exchange and clearing minimums, broker ticket fees, operational
overhead). Adding it would restore an interior optimal band width and remove the boundary
solution of Section~3 directly, rather than via the frequency cap used here as a proxy.

Two further caveats specific to this sample. The 22-day window is short -- one fifth of
the ES sample -- so the width-matched percentages should be read as indicative, and the
single $+3.22\%$ cell is likely optimistic. And the tape draws sides i.i.d.\ 50/50, so
option arrivals carry no predictable structure; real client flow is clustered, which would
change the balance between waiting for offsetting flow and hedging.

\section{Possible Improvements}

\begin{enumerate}
\item \textbf{Fixed per-ticket cost.} Replace the frequency cap with an explicit
  $C_{\text{ticket}}$ per child order. This is the single highest-value change: it makes
  the optimal band width interior and the static-vs-gamma comparison decision-relevant
  without an externally imposed constraint.
\item \textbf{Extend the sample.} 22 days is thin for distinguishing effects of 1--3\%.
  Extending the SPCX window, or bootstrapping across days, would put error bars on the
  width-matched edge.
\item \textbf{Volatility-conditioned band.} Gamma measures how fast delta moves per
  dollar of underlying; a realized-volatility divisor
  $1 + c_\sigma\hat\sigma/\tilde\sigma$ captures how many dollars the stock actually
  moves. Same plumbing, and on a name whose implied volatility spans 70--126\%, the
  intraday variation is likely larger than the factor of $2.1$ gamma offers.
\item \textbf{Calibrate $Y$ and the closing-cross haircut.} The square-root prefactor
  ($Y = 0.5$) and the $0.30$ auction haircut are assumptions, and the cost levels here are
  large enough that they deserve TCA support.
\item \textbf{Historical VWAP weights.} VWAP currently uses same-day realized volumes
  (mild look-ahead); a trailing time-of-day volume profile would make it implementable.
\item \textbf{Cross-name validation.} The intraday gamma range on SPCX, a factor of
  about $2.1$, bounds the achievable edge. Testing a more convex, event-driven book would show
  whether gamma conditioning scales with the gamma signal, as the mechanism predicts.
\end{enumerate}

\end{document}
